% year: תשע"ח
% type: מבחן
% semester: א
% moed: א
% number: 6ב
% points: 10
% categories: integral-applications, polar-parametric
% official_solution: yes
% source: exams/מבחנים/תשעח/מועד א תשעח.pdf

\begin{question}
(10 נק') נגדיר $r=f(\theta)=\theta^2-1$ בקואורדינטות קטביות. חשבו את אורך העקומה הזאת בין $\theta=1$ ל-$\theta=9$.

(במבחן מצורף איור של העקומה — ספירלה — על רקע רשת קטבית.)
\end{question}

\begin{hint}
אורך עקומה בקואורדינטות קטביות הוא $\int_\alpha^\beta\sqrt{r^2+\left(\frac{dr}{d\theta}\right)^2}\,d\theta$. בדקו שהביטוי מתחת לשורש הוא ריבוע שלם.
\end{hint}

\begin{solution}
אורך עקומה בקואורדינטות קטביות $r=f(\theta)$, $\alpha\le\theta\le\beta$:
\[ L=\int_\alpha^\beta\sqrt{r^2+\left(\frac{dr}{d\theta}\right)^2}\,d\theta. \]
כאן $\frac{dr}{d\theta}=2\theta$, ולכן
\[ r^2+\left(\frac{dr}{d\theta}\right)^2=(\theta^2-1)^2+4\theta^2=\theta^4-2\theta^2+1+4\theta^2=\theta^4+2\theta^2+1=(\theta^2+1)^2. \]
מכיוון ש-$\theta^2+1>0$, $\sqrt{(\theta^2+1)^2}=\theta^2+1$, ו-
\[ L=\int_1^9(\theta^2+1)\,d\theta=\left[\frac{\theta^3}{3}+\theta\right]_1^9=\left(243+9\right)-\left(\frac13+1\right)=252-\frac43=\frac{752}{3}. \]
\[ \boxed{L=\frac{752}{3}=250\tfrac23} \]
\textbf{הערה:} בפתרון הרשמי נכתב $252-\frac13+1=252\frac23$ — טעות סימן בהצבת הגבול התחתון; הערך הנכון הוא $252-\left(\frac13+1\right)=250\frac23$.
\end{solution}
